\documentclass[sigconf, screen, review, anonymous]{acmart}
%\documentclass[sigconf]{acmart}
%%
%% \BibTeX command to typeset BibTeX logo in the docs
\AtBeginDocument{%
  \providecommand\BibTeX{{%
    Bib\TeX}}}

%% Rights management information.  This information is sent to you
%% when you complete the rights form.  These commands have SAMPLE
%% values in them; it is your responsibility as an author to replace
%% the commands and values with those provided to you when you
%% complete the rights form.
\setcopyright{acmlicensed}
\copyrightyear{2018}
\acmYear{2018}
\acmDOI{XXXXXXX.XXXXXXX}
%% These commands are for a PROCEEDINGS abstract or paper.
\acmConference[Conference acronym 'XX]{Make sure to enter the correct
  conference title from your rights confirmation email}{June 03--05,
  2018}{Woodstock, NY}
%%
%%  Uncomment \acmBooktitle if the title of the proceedings is different
%%  from ``Proceedings of ...''!
%%
%%\acmBooktitle{Woodstock '18: ACM Symposium on Neural Gaze Detection,
%%  June 03--05, 2018, Woodstock, NY}
\acmISBN{978-1-4503-XXXX-X/2018/06}


%%
%% Submission ID.
%% Use this when submitting an article to a sponsored event. You'll
%% receive a unique submission ID from the organizers
%% of the event, and this ID should be used as the parameter to this command.
%%\acmSubmissionID{123-A56-BU3}

%%
%% For managing citations, it is recommended to use bibliography
%% files in BibTeX format.
%%
%% You can then either use BibTeX with the ACM-Reference-Format style,
%% or BibLaTeX with the acmnumeric or acmauthoryear sytles, that include
%% support for advanced citation of software artefact from the
%% biblatex-software package, also separately available on CTAN.
%%
%% Look at the sample-*-biblatex.tex files for templates showcasing
%% the biblatex styles.
%%

%%
%% The majority of ACM publications use numbered citations and
%% references.  The command \citestyle{authoryear} switches to the
%% "author year" style.
%%
%% If you are preparing content for an event
%% sponsored by ACM SIGGRAPH, you must use the "author year" style of
%% citations and references.
%% Uncommenting
%% the next command will enable that style.
%%\citestyle{acmauthoryear}

% Hashimoto added
\sloppy % 無理に詰めてoverfullになるのを防ぐ。
\usepackage{xcolor}
\newcommand{\TODO}[1]{\textcolor{red}{\textbf{TODO: #1}}}
\usepackage{booktabs}
\usepackage{comment}

% メモ: 致命的ではないので放置で良いが、vspaceとかのerrorがでているのはadjustwidthのせいかも??
\usepackage{changepage}
\usepackage{multirow} 
%%
%% end of the preamble, start of the body of the document source.
\begin{document}

%%
%% The "title" command has an optional parameter,
%% allowing the author to define a "short title" to be used in page headers.
\title{Manga109Script: Manga Script Dataset for Script2Layout}

%%
%% The "author" command and its associated commands are used to define
%% the authors and their affiliations.
%% Of note is the shared affiliation of the first two authors, and the
%% "authornote" and "authornotemark" commands
%% used to denote shared contribution to the research.
\author{Kengo Watanabe}
\email{kng.w@keio.jp}
\author{Shun Kitagawa}
\author{Yuro Okada}
\author{Michimasa Inaba}
\author{Atsushi Hashimoto}
\author{Satoshi Kurihara}


% 著者の書き方の参考
% \authornotemark[1]
% \affiliation{%
%   \institution{Institute for Clarity in Documentation}
%   \city{Dublin}
%   \state{Ohio}
%   \country{USA}
% }


% \author{Valerie B\'eranger}
% \affiliation{%
%   \institution{Inria Paris-Rocquencourt}
%   \city{Rocquencourt}
%   \country{France}
% }

% \author{Aparna Patel}
% \affiliation{%
%  \institution{Rajiv Gandhi University}
%  \city{Doimukh}
%  \state{Arunachal Pradesh}
%  \country{India}}

% \author{Huifen Chan}
% \affiliation{%
%   \institution{Tsinghua University}
%   \city{Haidian Qu}
%   \state{Beijing Shi}
%   \country{China}}

% \author{Charles Palmer}
% \affiliation{%
%   \institution{Palmer Research Laboratories}
%   \city{San Antonio}
%   \state{Texas}
%   \country{USA}}
% \email{cpalmer@prl.com}

% \author{John Smith}
% \affiliation{%
%   \institution{The Th{\o}rv{\"a}ld Group}
%   \city{Hekla}
%   \country{Iceland}}
% \email{jsmith@affiliation.org}

% \author{Julius P. Kumquat}
% \affiliation{%
%   \institution{The Kumquat Consortium}
%   \city{New York}
%   \country{USA}}
% \email{jpkumquat@consortium.net}

%%
%% By default, the full list of authors will be used in the page
%% headers. Often, this list is too long, and will overlap
%% other information printed in the page headers. This command allows
%% the author to define a more concise list
%% of authors' names for this purpose.
\renewcommand{\shortauthors}{Trovato et al.}

%%
%% The abstract is a short summary of the work to be presented in the
%% article.
\begin{abstract}
  In creative endeavors, a well-designed layout plays a crucial role. An effective layout enhances visual appeal and draws the reader's attention. Although this principle applies to various media, manga has a unique layout that reflects its narrative structure. Manga layouts are established in preliminary drafts called "Name," based on scripts and storytelling. Although layout generation models have evolved in domains such as Web UIs or Posters where visual clarity is prioritized, script-based layout generation (Script2Layout) has not been explored due to data scarcity. In this study, we created script data for approximately 20,000 pages from the Manga109 dataset and built the Manga109Script dataset linking scripts to manga layouts. Using this dataset, we developed manga layout generation models that directly take scripts as input using several LLMs as base models. Through comprehensive evaluation, we demonstrated that incorporating script information improves layout generation quality compared to existing methods, thereby establishing an effective baseline for the Script2Layout task.
\end{abstract}

%%
%% The code below is generated by the tool at http://dl.acm.org/ccs.cfm.
%% Please copy and paste the code instead of the example below.
%%
% \begin{CCSXML}
% <ccs2012>
%  <concept>
%   <concept_id>00000000.0000000.0000000</concept_id>
%   <concept_desc>Do Not Use This Code, Generate the Correct Terms for Your Paper</concept_desc>
%   <concept_significance>500</concept_significance>
%  </concept>
%  <concept>
%   <concept_id>00000000.00000000.00000000</concept_id>
%   <concept_desc>Do Not Use This Code, Generate the Correct Terms for Your Paper</concept_desc>
%   <concept_significance>300</concept_significance>
%  </concept>
%  <concept>
%   <concept_id>00000000.00000000.00000000</concept_id>
%   <concept_desc>Do Not Use This Code, Generate the Correct Terms for Your Paper</concept_desc>
%   <concept_significance>100</concept_significance>
%  </concept>
%  <concept>
%   <concept_id>00000000.00000000.00000000</concept_id>
%   <concept_desc>Do Not Use This Code, Generate the Correct Terms for Your Paper</concept_desc>
%   <concept_significance>100</concept_significance>
%  </concept>
% </ccs2012>
% \end{CCSXML}

% \ccsdesc[500]{Do Not Use This Code~Generate the Correct Terms for Your Paper}
% \ccsdesc[300]{Do Not Use This Code~Generate the Correct Terms for Your Paper}
% \ccsdesc{Do Not Use This Code~Generate the Correct Terms for Your Paper}
% \ccsdesc[100]{Do Not Use This Code~Generate the Correct Terms for Your Paper}
\begin{CCSXML}
<ccs2012>
<concept>
<concept_id>10010147.10010178.10010179</concept_id>
<concept_desc>Computing methodologies~Natural language processing</concept_desc>
<concept_significance>500</concept_significance>
</concept>
<concept>
<concept_id>10010147.10010178.10010224</concept_id>
<concept_desc>Computing methodologies~Computer vision</concept_desc>
<concept_significance>500</concept_significance>
</concept>
</ccs2012>
\end{CCSXML}

\ccsdesc[500]{Computing methodologies~Natural language processing}
\ccsdesc[500]{Computing methodologies~Computer vision}

%%
%% Keywords. The author(s) should pick words that accurately describe
%% the work being presented. Separate the keywords with commas.
\keywords{Layout Generation, Creative AI, Manga, Visual Storytelling, Script2Layout}
%% A "teaser" image appears between the author and affiliation
%% information and the body of the document, and typically spans the
%% page.
%  eye catch
\begin{teaserfigure}
  \includegraphics[width=\textwidth]{img/eye_catch_ACMMM.pdf}
  \caption{
Overview of the Script2Layout task. From left to right: (1) a manga page from the Manga109 dataset; (2) its corresponding script, generated using a vision large language model (VLLM) with visual prompting and speaker-aware dialogue sequencing; and (3) the layout predicted by a Script2Layout model based on the generated script. The script follows standard manga production conventions and includes scene headings, action descriptions, and dialogue. This setup allows layout models to make use of narrative structure for generating more coherent and expressive layouts. (Source: *Hanzai Kousyounin Minegishi Eitarou* by Takashi Kii)
}
  \Description{}
  \label{fig:eyecatch}
\end{teaserfigure}

\received{20 February 2007}
\received[revised]{12 March 2009}
\received[accepted]{5 June 2009}

%%
%% This command processes the author and affiliation and title
%% information and builds the first part of the formatted document.
\maketitle

\section{Introduction}
Storytelling takes many forms, from text-focused novels to visually-driven animation. Among these, manga has developed a unique expressive form that divides a flat surface into panels containing characters and dialogue. In typical manga production, this layout decision is first made in a draft called a "name." Based on scripts and plots, the layout, including panel size and character positioning, is crafted to enhance the presentation of the story, playing a vital role in manga appeal~\cite{fukaya_mangalayout}.

Layout generation models have primarily evolved targeting functional layouts such as Web UIs~\cite{Kikuchi2021_webUI_layout} and posters~\cite{tanaka_scipostlayout_2024, advertising_posters}, focusing on accurate information transmission~\cite{shi_intelligent_2023}. In contrast, layouts that combine storytelling and dramatic effects, like manga, are under-researched.
The absence of data linking stories to layouts is a key reason for this gap.
Script data are often produced internally and rarely released publicly, making it difficult to obtain. 

In this study, we utilized Vision Large Language Models (VLLMs) to create Manga109Script, a script dataset for approximately 20,000 manga pages based on the Manga109 dataset ~\cite{matsui2017sketch,aizawa_building_2020}\footnote{\url{http://www.manga109.org/en}}.
Recent advances in VLLMs have enabled high-precision extraction of textual information from images~\cite{zhang_vision-language_2024}. This capability has significantly contributed to research in Visual Storytelling, which extracts narratives from image sequences~\cite{oliveira_story_2024, li_benchmark_2025}.
Despite advancements in VLLMs, understanding content across multiple panels remains challenging~\cite{ikuta_mangaub_2024,vivoli_comicap_2024}.
To support VLLM limitations, we applied Visual Prompting~\cite{shtedritski_what_2023} with the Manga109 dataset annotations and employed a rule-based algorithm designed for panel reading order~\cite{sachdeva_manga_2024}, to generate high-quality script data.

Our scripts follow standard script formatting conventions, including {\it scene headings} indicating location, {\it action descriptions} detailing character movements and scenery, and {\it dialogue} (Fig.~\ref{fig:eyecatch}).
We consulted with manga experts to confirm the format fits professional workflows.
The generated script's quality was also evaluated by humans.

Using this dataset, we developed manga layout generation models that take scripts as direct input.
We comprehensively evaluated these models against existing layout generation methods and assessed the impact of script input to establish a baseline for the Script2Layout task. Dataset, and code access will be provided upon publication.

Our contribution is three-fold:
\begin{enumerate}
    \item We introduce an annotation-assisted method to generate script data from manga images using VLLMs, validated by human evaluators.
    \item We publish the first professional script-format manga storytelling dataset, covering 20k pages.
    \item We establish robust baselines for the novel Script2Layout task, demonstrating for the first time how script influences manga layout.
\end{enumerate}

\section{Related Work}
\subsection{Layout Generation}

Layouts crucially impact viewers' comprehensionand perception of visual content. The pursuit of excellent layouts across various media has driven research in layout generation models \cite{shi_intelligent_2023,shi_understanding_2023,huang_survey_2023}. Previous layout generation research primarily targets layouts emphasizing element visibility, such as mobile UIs~\cite{deka_rico_2017,jiang_coarse--fine_2022}, posters~\cite{seol_posterllama_2024,yang_posterllava_2024,tanaka_scipostlayout_2024, advertising_posters}, and PDFs~\cite{zhong_publaynet_2019,jiang_coarse--fine_2022}. 

To address the diverse design needs in these visual domains, various constraint-based generation approaches have been developed based on practical use cases. These approaches include generation conditioned solely on element types \cite{lee2020neural,jiang_layoutformer_2023}, on both element types and size specifications \cite{jiang_layoutformer_2023,kong2022blt}, and relational constraints between elements \cite{lee2020neural,kikuchi_constrained_2021}.
There are also interactive variations, such as layout completion that infers missing elements from partial layouts \cite{gupta2021layouttransformer}, and layout refinement that improves existing layouts~\cite{rahman2021ruite}. 

The development of models under these various constraints has enabled more accurate layout generation.
Layout generation models have evolved from early GANs~\cite{li_layoutgan_2019} and VAEs~\cite{arroyo_variational_2021, jyothi2021layoutvaestochasticscenelayout} to Transformers~\cite{jiang_layoutformer_2023} and recent discrete diffusion models~\cite{inoue2023layoutdm, zhang_layoutdiffusion_2023}.

Despite these technical advancements, generating layout from narrative context is still unaerexplored. This challenge arises from both a lack of datasets and the needs to convert sentimental dynamics in narratives into specific visual layouts. Modern LLMs are now capable of interpreting complex contexts and nuances effectively~\cite{LLMunderstandingNarrative}. 
Leveraging these capabilities, we propose Script2Layout, a task to generate layouts directly from scripts. We evaluate our approach against state-of-the-art methods such as LayoutDM\cite{inoue2023layoutdm}, LayoutFormer++\cite{jiang_layoutformer_2023}, and LayoutPrompter\cite{lin_layoutprompter_2023}, thereby demonstrating both its effectiveness and innovative potential.

\subsection{Manga Understanding as Visual Storytelling}
In the script2layout task, having high-quality script data for manga images is essential.
Constructing such datasets requires advanced techniques to extract narrative information from images.

The challenge of generating and recognizing narratives from a sequence of images is called visual storytelling~\cite{huang_visual_2016}. This line of inquiry was initiated by the pioneering work of Huang~\cite{huang_visual_2016} and subsequently evolved through early approaches that employed CNN-based feature extraction in conjunction with Recurrent Neural Network (RNN) models to capture sequential continuity from image inputs~\cite{yu2017hierarchicallyattentivernnalbumsummarization}. Later advancements introduced hierarchical attention mechanisms~\cite{nahian_hierarchical_2019}, scene graph-based semantic structure modeling~\cite{wang_storytelling_scenegraphs2020}, and causal relationship extraction methods~\cite{song2024causalstorylocalcausalattention} to facilitate the derivation of meaningful narrative structures from images. These investigations extend beyond mere image recognition to focus on capturing inter-image relationships and ensuring overall narrative coherence.

More recently, the rapid performance improvements observed in VLLMs such as GPT-4o~\cite{wu_gpt-4o_2024} and Gemini~\cite{team_gemini_2024} have enhanced the precise extraction of detailed context from images, creating new oppotunities in Visual Storytelling~\cite{yang_storyllava_2025}.
On the other hand, applying VLLM to manga images reduces performance due to the gap by manga's distinctive visual language, which differs from natural images. Namely, manga images are primarily rendered in black and white, follow unique reading orders, and consist of panels in which time and space are discretely represented.
These distinctive features make manga understanding a complex task composed of multiple intertwined subtasks. As illustrated in the CoMix benchmark~\cite{vivoli_comicspap_2025}, these include panel segmentation, dialogue region ordering, and character re-identification across panels~\cite{sachdeva_manga_2024}. Among them, character re-identification is notably difficult: even GPT-4 achieves less than 50\% accuracy on this task\cite{vivoli_comicspap_2025}. 

In our study, we tackle the inherent structural challenges of manga image understanding, which includes multiple difficult subtasks, by providing pre-processed inputs such as ordered dialogue sequences with speaker attribution~\cite{li_manga109dialog_2024} and annotated character regions to support visual identification~\cite{shtedritski_what_2023}. This design allows the model to bypass low-level visual parsing and focus on generating high-quality narrative scripts, thereby leveraging its vision-language capabilities more effectively.

Narrative datasets like Manga109Story\cite{chen_manga_2024} provide prose narratives for individual manga.
Our primary aim is constructing a {\it script-format} text dataset. Our Manga109Script dataset is developed with professional creators so as to fit into the conventional manga production flow.


\section{Manga109Script}
Manga109Script is a script dataset corresponding to 19,555 manga pages from the Manga109 dataset \cite{matsui2017sketch,aizawa_building_2020}. Manga109 contains manga images from 109 commercial Japanese manga works along with layout information such as panel, character, and text positions, made available for academic research purposes with proper creator permissions. Manga109Script reformats the two-page spread layout information from Manga109 into single-page units, linking each page to its corresponding script data.

The script format follows a style commonly used in Japanese manga production. Each script consists of three main components:
(1) Scene Headings, marked with a leading “◯” to denote location,
(2) Action Descriptions, written with a three-character indent to describe characters’ actions and environmental details, and
(3) Dialogue, where the speaker’s name is followed by their speech enclosed in quotation marks.
These components are illustrated in the center of Figure~\ref{fig:eyecatch}.

We adopted this script format because manga has primarily developed within Japan, and all the manga titles for which scripts were generated in this study were originally created through Japanese production processes. This design choice ensures that the dataset aligns with actual manga creation practices.

\subsection{Data Construction Process}

To construct Manga109Script, we employed GPT-4o to generate page-level scripts from manga images in the Manga109 dataset. To support the model’s understanding of character identity and dialogue context, we introduced two key inputs: (1) visual prompting, where the face of each character in the image was enclosed with a uniquely colored circle~\cite{shtedritski_what_2023}, and (2) pre-extracted dialogue annotations, including the textual content of each speech balloon and its associated speaker's name~\cite{li_manga109dialog_2024}. The ordering of dialogues was determined in three steps: First, each dialogue balloon was assigned to the panel with which it had the highest Intersection over Union (IoU). Second, the reading order of panels on the page was determined using the algorithm proposed by Sachdeva et al.\cite{sachdeva_manga_2024}. Finally, within each panel, dialogues were ordered based on their proximity to the top-right corner of the panel’s bounding box, following previous works~\cite{vivoli_comicap_2024, sachdeva_manga_2024, hinami_towards_2021} .

To help the model recognize continuous narratives across sequential pages within the constraint of limited token length, we structured our interaction with GPT-4o into the following three components for each target page $i$:

\begin{itemize}
    \item \textbf{$\alpha$: Initial comprehension prompting} \\
    We provided the full dialogue from the episode that includes page~$i$ and simulated an affirmative response from the model to promote narrative understanding.

    \begin{adjustwidth}{0.5em}{0pt} % 控えめなインデントに
    \setlength{\tabcolsep}{2pt}
    \begin{tabular}{@{}r l@{}}
    \texttt{\textbf{We:}} & \texttt{Here is the dialogue for the entire} \\
                          & \texttt{episode. Did you understand the story?} \\
    \texttt{\textbf{GPT-4o:}} & \texttt{Yes.} \\
    \end{tabular}
    \end{adjustwidth}

    \item \textbf{$\beta$($i$): Script generation for the $i$-th page} \\
    We provided GPT-4o with the processed image of page~$i$ (with visual prompting), the dialogue on page~$i$, and the summary generated in step $c(i{-}1)$, covering the story up to page~$i{-}1$. The model was then instructed to generate the script for page~$i$ in the Japanese screenwriting format.

    \begin{adjustwidth}{0.5em}{0pt}
    \setlength{\tabcolsep}{2pt}
    \begin{tabular}{@{}r l@{}}
    \texttt{\textbf{User:}} & \texttt{Here is the processed image for page $i$,} \\
                          & \texttt{the dialogue on page $i$, and the summary} \\
                          & \texttt{from $\gamma(i{-}1)$. Please generate the script} \\
                          & \texttt{for this page.} \\
    \texttt{\textbf{Agent:}} & \texttt{(Generates script for page $i$)} \\
    \end{tabular}
    \end{adjustwidth}

    \item \textbf{$\gamma$($i$): Summary generation up to the $i$-th page} \\
    We then prompted the model to summarize the story up to page~$i$, based on the dialogue history including $b(i)$.

    \begin{adjustwidth}{0.5em}{0pt}
    \setlength{\tabcolsep}{2pt}
    \begin{tabular}{@{}r l@{}}
    \texttt{\textbf{Usesr:}} & \texttt{Based on our previous interactions,} \\
                          & \texttt{please summarize the story so far.} \\
    \texttt{\textbf{Agent:}} & \texttt{(Generates summary up to the $i$-th page)} \\
    \end{tabular}
    \end{adjustwidth}
\end{itemize}

To generate the script for page $i$, we executed the sequence:
\[
\alpha \rightarrow \beta(i{-}1) \rightarrow \beta(i)
\]
Here, the arrow ($\rightarrow$) indicates the chronological order in which each prompt was issued to GPT-4o. Note that for $i = 1$, the sequence becomes $\alpha \rightarrow \beta(1)$, and $\beta(1)$ does not include a summary input.

To generate the summary for the story up to page~$i$, we extended the sequence as follows:
\[
\alpha \rightarrow \beta(i{-}1) \rightarrow \beta(i) \rightarrow \gamma(i)
\]

By iteratively applying this series of interactions from $i = 1$ to the final page of each manga title, we constructed the full Manga109Script dataset. The dataset statistics are summarized in Table~\ref{tab:statistics}.



\begin{table}[t]
\setlength{\abovecaptionskip}{2pt}
\caption{
Statistics of Manga109Script. "Distinct character types" refers to the number of unique characters identified per page, while "unique speaker types" indicates the number of distinct speaker IDs per page.
}
\centering
\begin{tabular}{l r r}
\toprule
\textbf{Element} & \textbf{Total Count} & \textbf{Average per Page} \\
\midrule
\multicolumn{3}{l}{\textit{Layout Elements}} \\
\hspace{0.5em}Pages & 19,555 & \---\\
\hspace{0.5em}Panels & 103,850 & 5.31 \\
\hspace{0.5em}Character instances & 157,865 & 8.70 \\
\hspace{0.5em}Distinct character types & 60,812 & 3.11 \\
\hspace{0.5em}Dialogue balloons & 146,549 & 7.49 \\
\hspace{0.5em}Unique speaker types & 55,778 & 2.85 \\
\midrule
\multicolumn{3}{l}{\textit{Script Elements}} \\
\hspace{0.5em}Scene headings & 26,931 & 1.38 \\
\hspace{0.5em}Action descriptions & 138,890 & 7.04 \\
\hspace{0.5em}Dialogues & 134,901 & 7.10 \\
\hspace{0.5em}Japanese characters (text) & 6,703,063 & 342.78 \\
\bottomrule
\end{tabular}
\label{tab:statistics}
\end{table}

\subsection{Quality Verification}
%The generated scripts have a specialized format unlike short captions. Hence, standard image caption evaluation metrics were difficult to apply. 
Unlike general image captioning datasets, our scripts lack ground truth.
Additionally, the script format provides contextual coherence rather than word overlap.
Thus, we used human evaluation with specific metrics to verify script quality.

Evaluation items included: (1) {\it Scene Heading Recall}, determining whether locations depicted in the manga were correctly identified in script scene headings; (2) {\it Scene Heading Precision}, determining whether locations mentioned in script scene headings actually appeared in the manga; (3) {\it Action Description Precision}, determining whether action descriptions generated for a target page actually match that page when mixed with randomly selected action descriptions from other pages; (4) {\it Out-of-Context Action Description Detection Accuracy}, determining whether action descriptions from other pages could be correctly identified as not belonging to the target page; and (5) {\it Dialogue Recall}, determining whether dialogue appearing in the manga was accurately transcribed in the script.

We randomly selected approximately 10 consecutive pages from each work to maintain story continuity, totaling 1,000 pages. The 1,000 pages were divided into two sets of 500 pages each, with six evaluators each reviewing one set, resulting in three evaluations per page. These evaluators received fair compensation.

We calculated each evaluation metric by averaging correct judgments across all elements: Scene Heading Recall (0.793), Scene Heading Precision (0.838), Action Description Precision (0.664), Out-of-Context Action Description Detection Accuracy (0.910), and Dialogue Recall (0.988).

\section{Experiments}
\begin{figure*}[t]
  \includegraphics[width=\textwidth]{img/method.pdf}
  \caption{
Illustration of the manga layout generation task. The upper half shows the Script2Layout approach, where layout elements (e.g., character names, dialogue, panels) are extracted from a natural language script using rule-based processing, including visual and speaker-aware cues. These elements, along with the script, are input into a Script2Layout model to predict bounding box coordinates for layout generation. In contrast, the lower half depicts a traditional layout generation pipeline that requires explicit category labels (e.g., \texttt{Speech1}, \texttt{Chara2}) as input. While conventional methods rely on structured element types, Script2Layout enables direct use of narrative scripts, allowing layout generation to be conditioned on storytelling intent.
}
  \Description{}
  \label{fig:method}
\end{figure*}

The primary goal of our experiments is to evaluate the effectiveness of using scripts as input for manga layout generation and to establish a baseline for Script2Layout. To this end, using Manga109Script we compare Script2Layout models, which are LLMs fine-tuned to generate layout coordinates conditioned on script input, with several existing state-of-the-art layout generation methods~\cite{inoue2023layoutdm, jiang_layoutformer_2023, lin_layoutprompter_2023} under multiple constraint settings, including the presence or absence of script input. 

\subsection{Problem Formulation}
We formulate manga layout generation as the task of predicting spatial positions for a given list of element categories. A layout $l$ consisting of $N \in \mathbb{N}$ elements is represented as $l = \{(c_1, \mathbf{b}_1), \dots, (c_N, \mathbf{b}_N)\}$, where $c_i$ denotes the category of the $i$-th element and $\mathbf{b}_i$ is its corresponding bounding box. In our experiments, we consider three types of element categories: \textit{characters}, \textit{speech balloons}, and \textit{panels}. Each bounding box $\mathbf{b}_i$ is quantized on a $32 \times 32$ grid. Given the number of elements $N$ and their categories $\{c_i\}_{i=1}^N$, the task is to predict their spatial positions $\{\mathbf{b}_i\}_{i=1}^N$ on the manga page.

This formulation reflects real-world manga production workflows, where artists compose layouts (known as “name” in Japanese production) based on a predefined script that specifies the categories of elements to be placed on each page. This formulation also bridges a critical gap between traditional layout generation models, which cannot directly accept script input, and Script2Layout models, which are designed to condition layout prediction on narrative scripts.


Under this formulation, we first describe how existing layout generation methods are applied to this task. We then introduce our Script2Layout models, highlighting the differences in how they utilize script information.

\paragraph{Existing Methods.}
The input and output formats of conventional methods are illustrated in the bottom part of Figure~\ref{fig:method}. Conventional layout generation methods require element categories to be predefined. In our experiments, categories are based on the three main types (\textit{characters}, \textit{speech balloons}, and \textit{panels}) and are further distinguished on a per-page basis. Specifically, we assign category as $c_i \in \{\texttt{Chara}_i \mid i=1,\dots,25\} \cup \{\texttt{Speech}_i \mid i=0,\dots,25\} \cup \{\texttt{Panel}\}$. Here, $\texttt{Chara}_i$ denotes the $i$-th character instance on a page, $\texttt{Speech}_i$ represents a speech balloon spoken by $\texttt{Chara}_i$ (or narration when $i = 0$). The number 25 was chosen to provide a buffer, as the maximum number of distinct characters in a single page of Manga109Script is 22. Based on this formulation, we evaluate three representative layout generation methods: LayoutDM~\cite{inoue2023layoutdm}, LayoutFormer++~\cite{jiang_layoutformer_2023}, and LayoutPrompter~\cite{lin_layoutprompter_2023}. The first two models are trained on the manga layout task. LayoutPrompter, in contrast, is a prompt-based system that does not involve model training; we implemented it using the GPT-4o API. Although the system is capable of accepting script inputs, we did not provide script information in our experiments due to input length limitations. Consequently, LayoutPrompter was evaluated under the same conditions as other non-script-based methods.We adopt this generating conditioned on types (Gen-T), where the category of each element is known, as the base condition throughout this study. In addition, we evaluate the existing methods under two extended constraint settings: 

\begin{itemize}
  \item \textbf{Gen-TS (Generation conditioned on types and sizes):} Each element's size (i.e., width and height) is given in addition to its category.
  \item \textbf{Gen-R (Generation conditioned on relationships):} Pairwise spatial relations (e.g., "element A is above element B") between layout elements are provided as constraints.
\end{itemize}

\paragraph{Script2Layout Models.}
Under the above formulation, we use the following LLMs as base for Script2Layout models: T5~\cite{raffel2020exploring_t5}, Swallow-13B\footnote{\url{https://swallow-llm.github.io/swallow-llama.ja.html }~(cited 2025-04-11)}, Sarashina-13B\footnote{\url{https://huggingface.co/sbintuitions/sarashina1-13b}~(cited-2025-04-11)}, and DeepSeek-R1-Distill-Qwen-14B-Japanese~\cite{deepseek-ai_deepseek-r1_2025}, selected for their strong Japanese language capabilities and fine-tuning flexibility. The input and output formats of Script2Layout models are illustrated in the upper part of Figure~\ref{fig:method}. Unlike conventional methods that require predefined category labels, Script2Layout models can directly process linguistic content. They take as input the raw content of layout elements, such as character names and speech texts, extracted from the script. These contents are passed to the tokenizer of each LLM without assigning predefined category labels, allowing the model to interpret them directly based on its pretrained knowledge. The model then outputs a quantized bounding box $\mathbf{b}_i = ((x_1, y_1), (x_2, y_2))$ for each element. In addition, we design input formats that allow the model to incorporate the script itself. Specifically, we consider the following three variants:
\begin{itemize}
  \item \textbf{No Script:} No script information is included in the input.
  \item \textbf{Plain:} Raw script text is included as input.
  \item \textbf{Panel-Inserted:} Script input is augmented with explicit panel segmentation markers, indicating how the script content is divided across individual panels.
\end{itemize}
In addition, we further evaluated Script2Layout models under the Gen-TS setting, combining it with the \textbf{No Script} and \textbf{Plain} input conditions. This additional experiment aims to assess the extent to which script information contributes to layout prediction when the sizes of all elements are provided in advance.

\subsection{Dataset Split}

For training and evaluation, we filtered out 90 pages from Manga109Script that contained more than 50 layout elements, as such cases can cause combinatorial explosion in constraint-based models (particularly under Gen-R). The remaining 19,465 pages were used in our experiments.

We split the data by manga title to avoid title-level leakage across splits. The split maintains an 8:1:1 ratio, with 15,333 pages for training, 2,007 for validation, and 2,125 for testing. The training set was used to fine-tune LLMs for Script2Layout and to train conventional baseline models from randomly initialized parameters. The validation set was used for early stopping, and the test set for final performance evaluation.

\subsection{Evaluation Metrics}

We evaluate layout generation performance using two widely adopted metrics in the layout generation literature: Fréchet Inception Distance (FID) and mean Intersection over Union (mIoU). These metrics collectively assess both the fidelity and diversity of the generated layouts with respect to ground truth.

\paragraph{FID}
FID measures the similarity between distributions of real and generated layouts in a learned feature space. Following LayoutDM~\cite{inoue2023layoutdm}, we use FIDNetV3 as the feature extractor, which has been shown to be better suited for layout features than traditional image-based FID backbones. For Script2Layout, FID is computed between generated layouts and real layouts in the test set. In addition, as a reference, we also compute the FID between layouts in the validation and test sets and report it as “Real data.”

\paragraph{mIoU}
We compute mean Intersection over Union (mIoU) by first greedily matching generated and ground-truth elements based on maximum IoU per category. This approach, also used in prior works~\cite{inoue2023layoutdm}, accounts for category-specific spatial accuracy and is well-suited to our task setup where the element categories are known or inferred. For each test sample, the mIoU is calculated as the average over all matched pairs of elements, and the final score is averaged over the entire test set.

These two metrics complement each other: FID captures global distributional similarity and diversity, while mIoU reflects instance-level alignment between predicted and ground-truth layouts.

\paragraph{Qualitative Evaluation.}
In addition to quantitative metrics, we also conduct a qualitative comparison by visualizing generated layouts for selected pages in the test set. For each target page, layouts generated by different models are compared side-by-side against the ground-truth layout. This allows us to visually assess how well each model captures spatial structure and how effectively script information contributes to layout coherence and accuracy.

\subsection{Evaluations and Discussions}
\begin{table*}[t]
\renewcommand{\arraystretch}{1.0}
\caption{
FID and mIoU scores under the Gen-T setting. Script2Layout models are evaluated with three types of script input. Most models outperform existing baselines, especially in mIoU, demonstrating the effectiveness of script-based guidance.
}
\label{tab:results}
\begin{tabular}{llrrrrrr}
\toprule
&\multicolumn{1}{r}{Script Type} & \multicolumn{2}{c}{No Script} & \multicolumn{2}{c}{Plain} & \multicolumn{2}{c}{Panel-Inserted} \\
\cmidrule(l{9em}){2-8}
Model & \multicolumn{1}{r}{Metric} 
 & \multicolumn{1}{c}{FID $\downarrow$} 
 & \multicolumn{1}{c}{mIoU $\uparrow$}
 & \multicolumn{1}{c}{FID $\downarrow$}
 & \multicolumn{1}{c}{mIoU $\uparrow$}
 & \multicolumn{1}{c}{FID $\downarrow$}
 & \multicolumn{1}{c}{mIoU $\uparrow$} \\
\midrule
\multicolumn{7}{l}{\it Layout Generation Models} \\%[-0.8ex]
\multicolumn{2}{l}{\hspace{0.5em}LayoutDM} & \underline{22.428} & 0.161 & - & - & - & - \\
\multicolumn{2}{l}{\hspace{0.5em}LayoutFormer++} & 28.402 & 0.140 & - & - & - & - \\
\multicolumn{2}{l}{\hspace{0.5em}LayoutPrompter} & 133.502 & 0.138 & - & - & - & - \\
\midrule[0.01em]
\multicolumn{7}{l}{\it Script2Layout Models} \\%[-0.8ex]
\multicolumn{2}{l}{\hspace{0.5em}T5 (Fine-Tuned)} & 29.549 & \textbf{0.201} & 28.744 & \textbf{0.210} & 31.287 & \underline{0.204} \\
\multicolumn{2}{l}{\hspace{0.5em}Swallow-13B (Fine-Tuned)} & 35.845 & 0.156 & \textbf{18.415} & 0.182 & \underline{21.557} & 0.184 \\
\multicolumn{2}{l}{\hspace{0.5em}Sarashina-13B (Fine-Tuned)} & \textbf{20.377} & \underline{0.192} & 24.191 & \underline{0.207} & 25.127 & \textbf{0.213} \\
\multicolumn{2}{l}{\hspace{0.5em}DeepSeek-R1-Japanese-14B (Fine-Tuned)} & 24.618 & 0.174 & \underline{22.502} & 0.186 & \textbf{21.094} & 0.184 \\
\midrule[0.01em]
\multicolumn{2}{l}{Real data} & 4.780 & & 4.780 & & 4.780 & \\
\bottomrule
\end{tabular}
\end{table*}

\paragraph{Quantitative Results under Gen-T with Varying Script Inputs.}

\begin{figure*}[t]
  \includegraphics[width=\textwidth]{img/generated_layouts_2.pdf}
  \caption{
Qualitative comparison of generated layouts. Each row shows one example: the real manga layout (left), followed by predictions from existing methods and Script2Layout models with different script input settings (No Script, Plain, Panel-Inserted). Script2Layout models generally produce layouts closer to the real ones, particularly when script information is provided. (Source: *Akuhamu* by Satoshi Arai, *Doll Gun* by Ryusei Deguchi, *Touta Mairimasu* by Shinji Saijo)
}
  \Description{}
  \label{fig:quantitative_results}
\end{figure*}

Table~\ref{tab:results} shows the FID and mIoU results for each model under the base constraint (Gen-T) across three script input conditions: No Script, Plain, and Panel-Inserted.

The FID results reveal several noteworthy trends. The best overall FID score (18.415) was achieved by Swallow-13B under the Plain condition, despite the model performing poorly in the No Script setting (FID = 35.845). This drastic improvement underscores the utility of incorporating script information: even models that underperform without scripts can significantly benefit from structured narrative input. Similarly, T5 and DeepSeek-R1 also demonstrated consistent improvements in FID when transitioning from No Script to Plain, indicating that three out of the four Script2Layout models benefitted from direct script input in terms of distributional similarity. These results collectively support the effectiveness of Plain script input for layout generation.

In contrast, Sarashina-13B achieved the best FID score among all models in the No Script setting (FID = 20.377).
While surpassing all existing baseline models in the No Script condition, the score slightly decreased with script input (FID = 24.191 for Plain and 25.127 for Panel-Inserted). This suggests that unstructured narrative input might introduce noise for the model, complicating its effectivity with raw content.
%This may indicate that the added complexity of processing unstructured narrative input could introduce noise for models already performing well based on raw element content alone.

The Panel-Inserted condition yielded improved FID only for DeepSeek-R1, which attained its best score under this setting (FID = 21.094). For the other models, this condition led to performance drops compared to Plain. This suggests that panel-wise segmentation cues, while potentially informative, may also introduce structural complexity that only models with strong reasoning capabilities, like DeepSeek-R1, are able to leverage effectively.

Among the existing layout generation models, LayoutDM achieved the best FID score (22.428). It is worth noting that LayoutDM has previously demonstrated strong performance on PubLayNet~\cite{zhong_publaynet_2019}, achieving an FID of 7.95 where the real data baseline is 6.25~\cite{inoue2023layoutdm}. In our setting, however, the increased FID (22.428) reflects the unique difficulty of the manga layout generation task. Notably, four model-condition pairs surpassed LayoutDM’s FID in our experiments: Sarashina-13B (No Script), Swallow-13B (Plain, Panel-Inserted), and DeepSeek-R1 (Panel-Inserted). Three of these four cases involved direct use of script input, further reinforcing the advantage of leveraging narrative context.

As for mIoU, all four Script2Layout models showed improved scores when script input was included, compared to their respective No Script counterparts. This demonstrates that script information contributes to more accurate placement of individual elements, especially in terms of spatial relationships and character positioning. Furthermore, with the exception of Swallow-13B under the No Script condition, all Script2Layout models outperformed the existing methods on mIoU. These results suggest that operating directly on the semantic content of layout elements—rather than relying solely on predefined categories—offers a tangible advantage in capturing fine-grained layout structure.

Taken together, these results highlight the value of incorporating script input into the layout generation process. Plain script input improves performance across models, especially in mIoU metrics. Script2Layout models, particularly when enhanced with script information, exhibit a advantage over conventional baselines.

\paragraph{Qualitative Analysis: Visualizing the Impact of Script Input.}

Figure~\ref{fig:quantitative_results} visualizes the layouts generated by each model for the same set of manga pages. These examples offer qualitative insights into how different models, especially Script2Layout models, respond to the inclusion of script input.

The top row illustrates a classic 4-panel manga (yonkoma). Notably, both DeepSeek and T5 produced layouts under the Panel-Inserted condition that closely resemble the structure of the original layout. This suggests that providing panel-wise segmentation cues helps the model understand the intended visual rhythm of 4-panel comics. A similar trend can be seen in Swallow: although its No Script output appears noisy and unstructured, its output under the Plain condition exhibits clearer four-panel alignment, reinforcing the value of script input for regularizing layout structure. Among the existing methods, LayoutPrompter produced surprisingly good results in this case. This may be attributed to its prompt-based design, which inherently biases the layout towards grid-like regularity. However, this same tendency leads to misaligned results in the middle and bottom examples, where strict symmetry is not appropriate.

In the second row, the target page includes a dramatic close-up of a character pointing a gun at his own temple. Across multiple Script2Layout models that utilize script input, we observe that large character regions are consistently placed at the top-left—mirroring the structure of the original page. This alignment can be attributed to textual cues in the script such as: “Norman stands frozen in shock. His eyes are wide open, with sweat beading on his forehead.” These narrative descriptions likely guided the models to allocate significant spatial prominence to the main character.

The bottom row presents a page with few speech balloons and a single dominant character holding a bamboo sword. Here again, Script2Layout models that incorporate script input generally place a tall character region centrally, with a speech balloon positioned to the upper right—mimicking the original composition. The script line “Ayumi Maniwa holds her bamboo sword at the ready, standing with a dignified posture.” likely encouraged such semantically meaningful layouts. This example demonstrates the models’ ability to ground spatial reasoning in both visual and narrative salience when provided with rich script input.

Overall, these results qualitatively affirm that script information, especially when formatted as panel-wise or plain narrative input, can substantially enhance layout plausibility, expressiveness, and alignment with storytelling intent.

\begin{comment}
\begin{table}[t]
\renewcommand{\arraystretch}{1.2} % 行間の調整（中央寄せをきれいに）
\centering
\caption{
Evaluation results of Existing layout generation methods under different constraint settings.
}
\begin{tabular}{l l r r}
\hline
Constraint & Models     & FID $\downarrow$& mIoU $\uparrow$ \\
\hline
\multirow{3}{*}{Gen-T} 
    & LayoutDM             & \textbf{22.428}          & 0.161 \\
    & LayoutFormer++       & \underline{28.402}          & 0.140 \\
    & LayoutPrompter       & 133.503         & 0.138 \\
\hline
\multirow{3}{*}{Gen-TS} 
    & LayoutDM             & \underline{20.514} & \textbf{0.271} \\
    & LayoutFormer++       & \textbf{17.487}    & \underline{0.208} \\
    & LayoutPrompter       & 58.993             & 0.148 \\
\hline
\multirow{3}{*}{Gen-R} 
    & LayoutDM             & \textbf{21.405}          & 0.053 \\
    & LayoutFormer++       & \underline{26.848}          & \textbf{0.158} \\
    & LayoutPrompter       & 130.502         & \underline{0.097} \\
\hline
& Real Data & 4.780 &  \\
\hline
\end{tabular}
\label{tab:constraint_existing_method}
\end{table}
\end{comment}
\begin{table}[t]
%\renewcommand{\arraystretch}{1.2}
\centering
\caption{
Evaluation results of existing layout generation methods under different constraint settings.
}
\begin{tabular}{l l r r}
\toprule
\multicolumn{1}{c}{Model} & \multicolumn{1}{c}{Constraint} & \multicolumn{1}{c}{FID $\downarrow$}& \multicolumn{1}{c}{mIoU $\uparrow$}\\
\midrule
\multirow{3}{*}{LayoutDM} 
    & Gen-T  &22.428 &  \underline{0.161} \\
    & Gen-TS & \textbf{20.514}    & \textbf{0.271} \\
    & Gen-R  &  \underline{21.405}             & 0.053 \\
\midrule
\multirow{3}{*}{LayoutFormer++} 
    & Gen-T  & 28.402             & 0.140 \\
    & Gen-TS & \textbf{17.487}    & \textbf{0.208} \\
    & Gen-R  & \underline{26.848} & \underline{0.158} \\
\midrule
\multirow{3}{*}{LayoutPrompter} 
    & Gen-T  & 133.503            & \underline{0.138} \\
    & Gen-TS & \textbf{58.993}    & \textbf{0.148} \\
    & Gen-R  & \underline{130.502}& 0.097 \\
%\midrule
%Real Data & & 4.780 &  \\
\bottomrule
\end{tabular}
\label{tab:constraint_existing_method}
\end{table}

\paragraph{Evaluating Conventional Methods under Size and Relational Constraints}
Table~\ref{tab:constraint_existing_method} reports the results of evaluating conventional layout generation methods under Gen-TS and Gen-R constraints in addition to the standard Gen-T setting. The Gen-R setting introduces explicit pairwise spatial relationships (e.g., “element A is above element B”) between layout elements. However, even with this additional relational information, the performance does not show significant improvement compared to the Gen-T baseline. Notably, none of the Gen-R results for existing methods approach the performance gains observed from using narrative scripts in Script2Layout models. This suggests that in the context of manga layout generation, script input offers a more effective and expressive form of structural guidance than manually defined pairwise spatial constraints.

In contrast, the Gen-TS setting, which supplements Gen-T with known element sizes, leads to substantial improvements in both FID and mIoU. For example, LayoutDM’s FID improves from 22.428 (Gen-T) to 20.514 (Gen-TS), and its mIoU rises significantly from 0.161 to 0.271. These results highlight the critical role that element size plays in layout generation performance. This trend is further supported by results from Script2Layout models, as presented in Table~\ref{tab:constraint_Script2Layoyut}.


\begin{table}[t]
\centering
\caption{
Performance of Script2Layout models under Gen-T and Gen-TS constraints, with and without script input.
%For brevity, model names are omitted; row order matches Table~\ref{tab:results} (T5, Swallow, Sarashina, DeepSeek).
}
\begin{comment}
\begin{tabular}{llrrrr}
\toprule
\hspace{10em}~&\multicolumn{1}{r}{Script Type} & \multicolumn{2}{c}{No Script} & \multicolumn{2}{c}{Plain} \\
\cmidrule(lr){2-6}
Constraint & Models & \multicolumn{1}{c}{FID $\downarrow$} & \multicolumn{1}{c}{mIoU $\uparrow$} & \multicolumn{1}{c}{FID $\downarrow$} & \multicolumn{1}{c}{mIoU $\uparrow$} \\
\hline
\multirow{4}{*}{Gen-T}  & T5                              & 29.549          & \textbf{0.201} & 28.744          & \textbf{0.210} \\
                        & Swallow                         & 35.845          & 0.156          & \textbf{18.415} & 0.182          \\
                        & Sarashina                       & \textbf{20.377} & \underline{0.192}    & 24.191          & \underline{0.207}    \\
                        & DeepSeek                        & \underline{24.618}    & 0.174          & \underline{22.502}    & 0.186          \\ \hline
\multirow{4}{*}{Gen-TS} & T5                              & \underline{19.445}    & \underline{0.379}    & \underline{19.734}          & \textbf{0.394} \\
                        & Swallow                         & 20.650          & 0.337          & 19.946    & 0.334          \\
                        & Sarashina                       & \textbf{18.165} & 0.358          & \textbf{18.783}          & 0.352          \\
                        & DeepSeek                        & 20.769          & \textbf{0.381} & 20.013 & \underline {0.383}    \\ \hline
                        & Real Data                       & 4.780           &                & 4.780           &                \\ \hline
\end{tabular}
\end{comment}
\begin{tabular}{llrrrr}
\toprule
&\multicolumn{1}{r}{Script Type} & \multicolumn{2}{c}{No Script} & \multicolumn{2}{c}{Plain} \\
\cmidrule{2-6}
\multicolumn{1}{c}{Model} & \multicolumn{1}{c}{Constraint} & \multicolumn{1}{c}{FID $\downarrow$} & \multicolumn{1}{c}{mIoU $\uparrow$} & \multicolumn{1}{c}{FID $\downarrow$} & \multicolumn{1}{c}{mIoU $\uparrow$} \\
\midrule
\multirow{2}{*}{T5} & Gen-T   & 29.549          & 0.201          & 28.744          & 0.210 \\
                    & Gen-TS  & \textbf{19.445} & \textbf{0.379} & \textbf{19.734} & \textbf{0.394} \\ 
\midrule
\multirow{2}{*}{Swallow} & Gen-T   & 35.845          & 0.156          & \textbf{18.415} & 0.182 \\
                         & Gen-TS  & \textbf{20.650} & \textbf{0.337} & 19.946          & \textbf{0.334} \\ 
\midrule
\multirow{2}{*}{Sarashina} & Gen-T   & 20.377          & 0.192 & 24.191          & 0.207 \\
                           & Gen-TS  & \textbf{18.165} & \textbf{0.358}          & \textbf{18.783} & \textbf{0.352} \\ 
\midrule
\multirow{2}{*}{DeepSeek} & Gen-T   & 24.618          & 0.174          & 22.502 & 0.186 \\
                          & Gen-TS  & \textbf{20.769} & \textbf{0.381} & \textbf{20.013}          & \textbf{0.383} \\
%\midrule
%                          & Real Data  & 4.780           &                & 4.780           &                \\ 
\bottomrule
\end{tabular}
\label{tab:constraint_Script2Layoyut}
\end{table}

\paragraph{Script2Layout under Gen-TS: Complementarity of Scripts and Size Constraints}
The Gen-TS results for Script2Layout models demonstrate consistent performance gains over Gen-T across all four models. For example, the T5 (No Script) improves from 0.201 to 0.379 in mIoU, and T5 (Plain) achieves the highest mIoU of 0.394. These findings affirm the importance of providing size information in layout generation, regardless of whether scripts are used.

When comparing the effect of script usage in the Gen-TS condition, we identified two models (DeepSeek, T5) that showed improvements in mIoU and another two models(DeepSeek, Swallow) that improved in FID. While the improvement in the Gen-TS condition is not as pronounced as in the Gen-T setting, our findings indicate that incorporating script content remains beneficial even when detailed size are provided. Essentially, semantic constraints (from scripts) offer valuable contextual understanding, whereas geometric constraints (like size) enhance structural precision. In summary, these results underscore the potential benefits of balancing and integrating both sources of information, which could be a promising direction for future research.

\section{Conclusion}

In this work, we introduced Script2Layout, a benchmark framework for manga layout generation using natural language scripts. To support this task, we constructed Manga109Script, a large-scale dataset of 19,555 page-level scripts aligned with Manga109.

We implemented Script2Layout using multiple fine-tuned LLMs and compared them with existing layout generation methods under different constraints, including Gen-T, Gen-TS, and Gen-R. Experiments using FID and mIoU metrics show that Script2Layout models consistently outperform baselines, especially in mIoU. Incorporating script input proved effective, and adding size constraints (Gen-TS) further improved performance. \\

These results highlight the complementary role of semantic and size information in improving layout generation. Future work includes modeling narrative continuity and integrating manga-specific elements like character posture.

%\section*{Acknowledgments}
%This research was supported by NEDO's Next-Generation AI Technology Development Project "Development of Interactive Story-based Content Creation Infrastructure."


%% consistent spelling of the heading.
%\begin{acks}
%\end{acks}

%%
%% The next two lines define the bibliography style to be used, and
%% the bibliography file.
%\newpage
%\balance % acmartが自動でやってくれるはずだが、warningが出ていたのでmanualで\balanceを挿入。
\bibliographystyle{ACM-Reference-Format}
\bibliography{script2layout}


%%
%% If your work has an appendix, this is the place to put it.
%\appendix

%\section{Research Methods}

\end{document}
\endinput
%%
%% End of file `sample-sigconf.tex'.
